\songtitle{Ventana de noviembre}

\tag*{2026}\tag{rock-en-espanol}\tag{spanish-lyrics}
\begin{notes}
Lyrics by SUNO based on a Gemini Notebook summary of chapter 3 (2001--2008) of Carlos Thompson's memoir \emph{A Life in World Events}.
\end{notes}
\begin{style}
Rock en español con pulso folk narrativo, tempo medio, guitarra acústica punteada y eléctrica al frente; verso íntimo y detallista, pre-coro con tensión creciente, coro amplio y cantable con coros en grupo; voz principal cercana con dobles en el estribillo, ecos cortos en remates, rasgueos abiertos, bombo seco, armónica y un puente que cae a voz y guitarra antes del último golpe, folk, rock, world
\end{style}
\begin{lyrics}
\songpart{Verse 1}
En Bogotá comenzó\\
la casa prestada y tibia,\\
con el techo de los padres\\
y la mesa ya servida.

No era el borde del vacío,\\
no era el cero por entrar,\\
pero debajo de esa calma\\
algo empezó a chocar.

La oficina fue una jaula,\\
el lenguaje, un muro gris,
él buscaba en otras redes\\
lo que el cubículo no da.

\songpart{Pre-Chorus}
Y cuando cayó el mundo\\
en una sola mañana,\\
los viejos medios tropezaban\\
y la red ya lo nombraba.

\songpart{Chorus}
Se abrió la ventana\\
se abrió la ventana\\
y entró la verdad.

Se abrió la ventana\\
se abrió la ventana\\
ya no pudo mirar atrás.

\songpart{Verse 2}
Luego vino otro mapa,\\
otra ciudad, otro metal,\\
trabajar frente a un nombre\\
que sonaba a capital.

Pero la rueda se partió\\
por la mitad en dos mil dos,\\
y entre cursos, redes y aulas\\
se le fue la voz.

En la universidad\\
se sentaba a ver pasar\\
las ventanas de los salones\\
sin atreverse a entrar.

\songpart{Pre-Chorus}
Y el tiempo, tan callado,\\
le fue dejando la cuenta,\\
como un perro en la esquina\\
mordiéndole la puerta.

\songpart{Chorus}
Se abrió la ventana\\
se abrió la ventana\\
y entró la verdad.

Se abrió la ventana\\
se abrió la ventana\\
ya no pudo mirar atrás.

\songpart{Bridge}
Hubo sangre y hubo duelo,\\
hubo un cuerpo que faltó,\\
hubo un hijo que llegaba\\
cuando el suelo se movió.

Y aun así, entre ruinas,\\
una clave floreció:\\
aprendió de la tormenta\\
a hacer red con el dolor.

\songpart{Verse 3}
Llegó el niño, llegó el aire,\\
y una placa se ganó,\\
mientras caían empresas\\
y el país se desarmó.

De pantalla en pantalla\\
fue juntando una nación,\\
de la página al relato,\\
de la marcha a la acción.

No quiso volver al ciclo\\
del empleado en su rincón,\\
prefirió mover la historia\\
con código y con voz.

\songpart{Chorus}
Se abrió la ventana\\
se abrió la ventana\\
y entró la verdad.

Se abrió la ventana\\
se abrió la ventana\\
ya no pudo mirar atrás.

\songpart{Outro}
Y una noche de noviembre\\
soltó al fin su señal,\\
un mensaje diminuto\\
para empezar a andar.

Se abrió la ventana\\
se abrió la ventana\\
y empezó a hablar.
\end{lyrics}

